% ═══════════════════════════════════════════════════════════════════
% ЧАСТЬ I (фон). МАТЕМАТИЧЕСКИЙ ФОН ПРОГРАММЫ
% ═══════════════════════════════════════════════════════════════════

\section{Математический фон программы}\label{sec:background}

Этот раздел вводит все математические понятия, которыми пользуется
программа, в объёме, достаточном для самостоятельного чтения
монографии. Каждое понятие определено, объяснено и связано с
конкретной ролью в конвейере. Читатель, знакомый с римановыми
поверхностями и структурами Ходжа, может перейти сразу к
разделу~\ref{sec:torus}.

\subsection{Римановы поверхности и род}

\emph{Риманова поверхность} --- одномерное комплексное
многообразие: пространство, локально устроенное как область
комплексной плоскости, с согласованными переходами голоморфных
координат. Простейшие примеры --- сфера Римана $\mathbb{P}^1$ (род
$0$), тор $\CC/\Lambda$ (род $1$). Род --- целочисленный
топологический инвариант: число «ручек» поверхности; алгебраически
$g=\dim H^0(C,\Omega^1)$ --- размерность пространства голоморфных
форм. Для гладкой плоской кривой степени $m$ род равен
$\tfrac{(m-1)(m-2)}{2}$ --- эта формула (Плюккера) даёт род
квартики Клейна $3$ и, в другом прочтении (лемма~\ref{lem:genus}),
род кривой Ферма.

\emph{Кривая Ферма} $C_N=\{x^N+y^N=z^N\}$ --- плоская кривая
степени $N$; её род $(N-1)(N-2)/2$ растёт квадратично: при $N=15$
это $91$, при $N=30$ --- $406$. Кривая обладает большой группой
симметрий $\mu_N\times\mu_N$: умножения координат корнями из
единицы. Именно эта группа организует весь дальнейший счёт ---
формы и циклы раскладываются по её характерам.

\subsection{Голоморфные формы и периоды}

\emph{Голоморфная 1-форма} --- дифференциал вида $\omega=f(z)\,dz$
с голоморфной $f$ в каждой карте. На кривой рода $g$ их пространство
$g$-мерно; на кривой Ферма базис выписан явно (лемма~\ref{lem:forms}):
$\eta_{i,j}=x^iy^j\,dx/y^{N-1}$.

\emph{Период} --- интеграл формы по замкнутому циклу. Циклы
образуют решётку $H_1(C,\ZZ)\cong\ZZ^{2g}$; выбор базиса циклов
превращает совокупность периодов в \emph{период-матрицу} размера
$g\times2g$. Период-матрица кодирует комплексную структуру кривой
восстановимо: кривая восстанавливается по своей якобиановой
структуре (теорема Торе). На кривой Ферма каждый период имеет
замкнутую форму $\frac1N\zeta^{ra+sb}\Omega_{a,b}$
(лемма~\ref{lem:closedform}) --- трансцендентный множитель
$\Omega_{a,b}$ (бета-интеграл) и алгебраическая фаза (корень
единицы) разделены; это разделение --- техническая основа всей
$\Ga$-нормировки.

\emph{Якобиан} $\Jac(C)$ --- комплексный тор $\CC^g/H_1(C,\ZZ)$,
снабжённый главной поляризацией; это абелево многообразие,
универсальный объект, в который кривая вкладывается через
интегрирование форм (отображение Абеля--Якобиана). Для квартики
Клейна якобиан изогенен кубу эллиптической кривой с CM-полем
$\QQ(\sqrt{-7})$ --- теорема Рибета, вокруг которой построен
стенд раздела~\ref{sec:klein}.

\subsection{Соотношения Римана и структуры Ходжа}

Период-матрица в симплектическом базисе обязана удовлетворять
\emph{соотношениям Римана}: симметрия $\Omega^T=\Omega$ и
положительная определённость $\im\Omega$. На кривой они доказываются
билинейным тождеством Римана (теорема~\ref{th:riemann}) ---
единственным аналитическим инструментом, который программа
использует в этом блоке. Пространство $H^1$ распадается на
\emph{структуру Ходжа} $H^1=H^{1,0}\oplus H^{0,1}$ --- собственные
пространства оператора $\sqrt{-1}$ на комплексификации; форма
пересечений задаёт \emph{поляризацию}: условия Римана--Ходжа
(теорема~\ref{th:polarization}). Гипотеза Ходжа --- утверждение о
том, что рациональные классы, ортогональные всем трансцендентным
направлениям, являются классами алгебраических циклов; программа
строит вычислимую инфраструктуру вокруг этого утверждения ---
поляризации, решётки, индексы --- и доводит каждую компоненту до
точной сертификации.

\subsection{Спиновые структуры и инвариант Арфа}

\emph{Спиновая структура} на кривой --- выбор квадратного корня из
канонического расслоения; на кривой рода $g$ их ровно $2^{2g}$.
Чётность структуры определяется \emph{инвариантом Арфа} ---
квадратичной формой на $H^1(C,\ZZ/2)$. Чётность управляет
специальной тета-константой: чётные структуры имеют нулевую
тета-константу. Счёт чётных/нечётных структур по родам ($1/3$ на
роде~1, $6/10$ на роде~2, $28/36$ на роде~3) --- классический
результат, воспроизведённый программой полным перечислением
(раздел~\ref{sec:e8}); сам счёт --- квадратичная форма Арфа на
$\ZZ/2^{2g}$, и перебор $64$ структур на роде~3 занимает
микросекунды.

\subsection{Циклотомические поля и корни единицы}

\emph{Корни единицы} $\zeta_N^k=e^{2\pi ik/N}$ образуют циклическую
группу $\mu_N$; поле $\QQ(\zeta_N)$ --- циклотомическое поле
степени $\phi(N)$. Фазовая решётка стенда $\frac1N\ZZ[\zeta_N]$ ---
модуль над кольцом целых циклотомического поля; знаменатель $N$ ---
это то, что теорема~\ref{th:normalization} фиксирует как
«дробные доли $\pi$». Функция Мёбиуса $\mu(d)$ ---
мультипликативный инвариант с значениями
$\{0,\pm1\}$; формула \eqref{eq:mobius} выражает через неё перепись
характеров по проводникам --- это классическая схема
включений-исключений, и её точность (целые числа!) делает перепись
сертификатом первого типа.

\subsection{Кривые, решётки и CM-поля}

Эллиптическая кривая с \emph{комплексным умножением} (CM) имеет
периодное отношение в мнимом квадратичном поле; например,
$\tau=(1+\sqrt{-7})/2$ для компоненты якобиана Клейна. Решётка
$\ZZ\tau+\ZZ$ имеет спектр лапласиана $|m\tau-n|^2$ ---
универсальную формулу сертификата~G1. Числа классов CM-полей
малых дискриминантов ($h(-7)=1$, $h(-3)=1$, $h(-15)=2$) делают
все упомянутые объекты вычислимыми в целых числах --- этим
программа систематически пользуется: каждый CM-объект возникает
вместе со своим точным минимальным многочленом
($4X^2-7$, $4X^3-1$ и т.д.), и никакая величина не остаётся
«просто числом».

\subsection{Метод tanh-sinh: почему он независим}

Протокол сверяет $\Ga$-формы прямым интегрированием методом
tanh-sinh (квадратура двойной экспоненты). Метод применяется к
\emph{гладким} подынтегральным функциям: замена $t=u^N/2$ на
половинах отрезка снимает концевые особенности бета-интеграла, и
в коде квадратуры не появляется ни одной $\Ga$-функции --- это
гарантирует независимость сертификата: ошибка в формуле
$\Ga(a/N)\Ga(b/N)/\Ga((a+b)/N)$ не может «поддержать» сама себя,
потому что эталон считается совсем другим алгоритмом. Согласие двух
методов до $10^{-36}$ при $dps=35$ --- сертификат корректности
формулы, а не просто численного совпадения.

\subsection{Нормальная форма Смита и индексы}

\emph{Нормальная форма Смита} (SNF) целочисленной матрицы ---
диагональ $d_1|d_2|\dots$ с точными целыми инвариантными
множителями; она вычисляется целочисленными элементарными
преобразованиями без округлений. Для поляризационных решёток SNF
даёт \emph{тип поляризации} --- целочисленный инвариант главного
поляризованного многообразия. Произведение инвариантных множителей
равно абсолютному значению определителя --- для стенда
$N=15/30$ это $1\cdot1\cdot5\cdot5\cdot15\cdot15\cdot15\cdot15
=1265625=1125^2$, согласование с дискриминантом --- целочисленный
сертификат (теорема~\ref{th:H1}). Индекс подрешётки ---
отношение определителей подрешётки и насыщения; алгоритм Эрмита--
Смита вычисляет базис насыщения точно, что делает индексы блоков
(теорема~\ref{th:indexcomp}) вычислимыми в строгом смысле.
